
Pre-training corpus curation is among the most resource-intensive components of large language model development. Decisions about which documents to include, how aggressively to deduplicate, and which high-quality sources to prioritize (e.g., academic papers, Wikipedia) are made under the assumption that higher-quality data produces models with stronger generalization --- not only higher average benchmark scores, but better balance between knowledge-intensive and commonsense reasoning capabilities. Despite the practical stakes of this assumption, it has rarely been tested under controlled conditions at matched training scales.

The specific hypothesis examined in this work is: \textit{at approximately 300 billion training tokens, a model trained on a more carefully curated corpus will achieve a higher MMLU/HellaSwag generalization balance ratio than a model trained on a less curated corpus.} This hypothesis is operationalized by comparing OLMo-7B \citep{Groeneveld2024OLMo}, trained on Dolma \citep{Soldaini2024Dolma} (multi-stage quality filtering, URL blocklisting, deduplication, and explicit inclusion of academic sources), against Pythia-6.9B \citep{Biderman2023Pythia}, trained on The Pile \citep{Gao2020Pile} (minimal curation, basic deduplication and language filtering). Both models provide publicly available intermediate checkpoints at approximately matched training token counts, enabling a controlled matched-scale comparison.

Prior work comparing curated and uncurated corpora has consistently confounded corpus quality with training duration, model architecture, or both. Falcon models trained on RefinedWeb \citep{Penedo2023RefinedWeb} differ from The Pile-trained models in architecture; Pythia-dedup vs. Pythia comparisons \citep{Biderman2023Pythia} measure deduplication in isolation. No prior study has evaluated MMLU/HellaSwag generalization balance ratios for an architecture-paired comparison at matched training token counts.

The primary finding is a direct refutation of the hypothesis: Pythia-6.9B achieves a higher MMLU/HellaSwag ratio (0.5650 vs. 0.5383) at matched training scale, and this advantage is consistent across all four benchmarks evaluated. The directional refutation is statistically robust (95\% CI entirely below zero, Cohen's d = $-2.732)$. More structurally informative is the observation that both models achieve identical HellaSwag accuracy (0.4580), irrespective of corpus curation quality, at this training scale. When the denominator of the generalization balance ratio is a constant, the ratio reduces to a noisy proxy for MMLU differences alone --- differences that cannot be attributed to corpus quality without resolving the architecture confound.

This paper makes four contributions:

\begin{enumerate}
\item \textbf{Matched-scale empirical evaluation.} The first systematic evaluation of MMLU/HellaSwag generalization balance as a corpus-quality discriminator, comparing Pythia-6.9B (The Pile, step143000 $\approx$ 299.9B tokens) and OLMo-7B (Dolma, step68000 $\approx$ 301B tokens) at controlled matched scale. The corpus-quality prediction is directly refuted (ratio difference = $-0.0267$, 95\% CI [$-0.0447, -0.0069]$, one-sided p = 0.996 for the test that OLMo $\geq$ Pythia).
\end{enumerate}

\begin{enumerate}
\item \textbf{Metric sensitivity analysis.} HellaSwag 0-shot commonsense accuracy converges to identical values (0.4580) for both models at approximately 300 billion tokens, demonstrating that the MMLU/HellaSwag ratio is not a reliable corpus curation quality discriminator at this training scale in a cross-architecture comparison.
\end{enumerate}

\begin{enumerate}
\item \textbf{Methodological characterization.} The conditions under which cross-architecture, single-checkpoint comparisons fail to support causal corpus-quality claims are characterized: when a ratio metric's denominator saturates, the ratio loses its intended discriminative property, and architecture differences in the numerator cannot be attributed to data quality.
\end{enumerate}

\begin{enumerate}
\item \textbf{Scope-qualified null result.} The negative finding is explicitly scoped to approximately 300 billion training tokens, the GPT-NeoX/LLaMA-style architecture pairing, and the MMLU/HellaSwag ratio metric. It does not rule out corpus quality effects at other training scales, with architecture-matched designs, or using different generalization metrics.
\end{enumerate}

